%; whizzy paragraph -pdf xpdf -latex ./whizzypdfptex.sh
%; whizzy-paragraph "^\\\\begin{frame}"
% latex beamer presentation.
% platex, latex-beamer でコンパイルすることを想定。

%     Tokyo Debian Meeting resources
%     Copyright (C) 2009 Junichi Uekawa
%     Copyright (C) 2009 Nobuhiro Iwamatsu

%     This program is free software; you can redistribute it and/or modify
%     it under the terms of the GNU General Public License as published by
%     the Free Software Foundation; either version 2 of the License, or
%     (at your option) any later version.

%     This program is distributed in the hope that it will be useful,
%     but WITHOUT ANY WARRANTY; without even the implied warreanty of
%     MERCHANTABILITY or FITNESS FOR A PARTICULAR PURPOSE.  See the
%     GNU General Public License for more details.

%     You should have received a copy of the GNU General Public License
%     along with this program; if not, write to the Free Software
%     Foundation, Inc., 51 Franklin St, Fifth Floor, Boston, MA  02110-1301 USA

\documentclass[cjk,dvipdfmx,12pt]{beamer}
\usetheme{Tokyo}
\usepackage{monthlypresentation}

%  preview (shell-command (concat "evince " (replace-regexp-in-string
%  "tex$" "pdf"(buffer-file-name)) "&"))
%  presentation (shell-command (concat "xpdf -fullscreen " (replace-regexp-in-string "tex$" "pdf"(buffer-file-name)) "&"))
%  presentation (shell-command (concat "evince " (replace-regexp-in-string "tex$" "pdf"(buffer-file-name)) "&"))

%http://www.naney.org/diki/dk/hyperref.html
%日本語EUC系環境の時
\AtBeginDvi{\special{pdf:tounicode EUC-UCS2}}
%シフトJIS系環境の時
%\AtBeginDvi{\special{pdf:tounicode 90ms-RKSJ-UCS2}}

\title{golangの''net/http/fcgi''を使ってREST APIアプリを作る}
\subtitle{Debian / Ubuntu ユーザーミートアップ in 札幌 2019.05}
\author{Norimitsu Sugimoto (杉本 典充) \\dictoss@live.jp}
\date{2019-05-31}
\logo{\includegraphics[width=8cm]{image200607/openlogo-light.eps}}

\begin{document}

\frame{\titlepage{}}


\section{アジェンダ}

\begin{frame}{アジェンダ}
  \begin{itemize}
  \item 自己紹介
  \item golangについて
  \item Debianとgolang
  \item goのアプリケーションを動かす
  \item おわりに
  \item 参考資料
  \end{itemize}
\end{frame}


\section{自己紹介}

\begin{frame}{自己紹介}
  \begin{itemize}
  \item Norimitsu Sugimoto (杉本 典充)
  \item dictoss@live.jp
  \item Twitter: @dictoss
  \item 大学生のときからDebianを使っています
  \item 仕事はソフトウェア開発者をやってます
  \end{itemize}
\end{frame}


\section{golangについて}
\emtext{golangについて}

\begin{frame}[containsverbatim]{golangについて}
  \begin{itemize}
  \item 2009/11/10にgoogleが発表したオープンソースのプログラミング言語
  \item 最新の安定版「go1.12」のリリースは 2019/2/25
  \item プログラミングの特徴
    \begin{itemize}
    \item 変数は静的な型付けを行う
    \item ネットワークやマルチスレッド処理を行う豊富な標準ライブラリを持つ
    \item structでデータ構造を定義でき、classはない（interfaceはある）
    \end{itemize}
  \item 特徴
    \begin{itemize}
    \item "go build" コマンドでOSネイティブな実行ファイルをビルドする
      \begin{itemize}
      \item Linux、macOS、Windows、FreeBSD 向けの実行ファイルをビルド可能
      \end{itemize}
    \end{itemize}
  \item golangで実装しているアプリケーション例
    \begin{itemize}
    \item \url{https://github.com/golang/go/wiki/GoUsers}
    \end{itemize}
  \end{itemize}
\end{frame}


\section{Debianとgolang}

\begin{frame}[containsverbatim]{Debianとgolang}
  \begin{itemize}
  \item Debian 9 Stretch
    \begin{itemize}
    \item golang-1.7.4
    \item golang-1.11.5（stretch-backports）
    \end{itemize}
  \item Debian 10 Buster
    \begin{itemize}
    \item golang-1.11.6
    \end{itemize}
  \end{itemize}
\end{frame}


\section{goのアプリケーションを動かす}
\emtext{goのアプリケーションを動かす}

\subsection{動作環境とデモアプリケーション}

\begin{frame}[containsverbatim]{動作環境とデモアプリケーション}
  \begin{itemize}
  \item Debian GNU/Linux 10 Buster amd64
  \item golang パッケージ（golang-1.11）のgoコマンドを利用する
  \item デモアプリケーション
    \begin{itemize}
    \item \url{https://github.com/dictoss/golang-practice/tree/master/src/github.com/dictoss/fcgi_server}
    \item FastCGIとして動作するREST APIサーバ
    \item このアプリケーションをデーモンとして実行
    \end{itemize}
  \item nginx
    \begin{itemize}
    \item http のリクエストを待ち受ける
    \item リクエストを FastCGI のバックエンドへ proxy する
    \end{itemize}
  \end{itemize}
\end{frame}


\subsection{パッケージのインストール}

\begin{frame}[containsverbatim]{パッケージのインストール}

aptコマンドで必要なパッケージをインストール

\begin{commandline}
# apt update
# apt install golang nginx curl git
\end{commandline}

goコマンドのバージョンを確認

\begin{commandline}
$ go version
go version go1.11.6 linux/amd64
\end{commandline}

\end{frame}


\subsection{ビルドと実行}

\begin{frame}[containsverbatim]{ビルドと実行}

ソースコードを取得

\begin{commandline}
$ git clone https://github.com/dictoss/golang-practice.git
\end{commandline}

ビルドして実行ファイルを生成

\begin{commandline}
$ cd golang-practice/src/github.com/dictoss/fcgi_server/
$ go build fcgi_server.go
$ file fcgi_server
fcgi_server: ELF 64-bit LSB executable, x86-64,
version 1 (SYSV), dynamically linked,
interpreter /lib64/ld-linux-x86-64.so.2,
Go BuildID=NgM3NZ0-K_rN-eMIo4T1/8VObv2x-PXhaCtETqR2t
/4HbNDkw2hyrGBQ03TFnI/OhB_Jj--KNJF6NYEIL6z,
not stripped
\end{commandline}

実行できるか確認

\begin{commandline}
$ ./fcgi_server
\end{commandline}

\end{frame}


\subsection{nginxのfastcgi設定}

\begin{frame}[containsverbatim]{nginxのfastcgi設定}

webサーバの /gofcgi のパスへアクセスされたときにFast CGIでproxyする

\begin{commandline}
# vi /etc/nginx/sites-available/default
server {
    location ~ ^/gofcgi/(.*)$ {
        include /etc/nginx/fastcgi_params;
        fastcgi_pass 127.0.0.1:9000;
    }
}
\end{commandline}

\begin{commandline}
# nginx -c /etc/nginx/nginx.conf
# systemctl restart nginx
\end{commandline}

\end{frame}


\subsection{動作確認}

\begin{frame}[containsverbatim]{動作確認}

動作確認

\begin{commandline}
  $ cat hello2_request.json | \
  curl -v http://localhost/gofcgi/json/hello2/ \
  -H 'Content-Type:application/json' --data @-

  {``req_msg'':''Nice to meet you.'',
   ''res_msg'':''Me too !!''}
  \end{commandline}
\end{frame}


\subsection{デーモン設定}

\begin{frame}[containsverbatim]{デーモン設定（1）}

systemdのserviceファイル

\begin{commandline}
$ cat gofcgi-apid.service
[Unit]
Description=REST API server for golang
After=network.target

[Service]
Type=simple
ExecStart=/usr/local/gofcgi-apid/fcgi_server
  -logpath ``/usr/local/gofcgi-apid/log.txt''
  -fcgi_listen_addr ``127.0.0.1:9000''
  -fcgi_url_prefix ``/gofcgi''
ExecStop=/bin/kill -HUP $MAINPID
TimeoutStopSec=5
KillMode=mixed

[Install]
WantedBy=multi-user.target
\end{commandline}

\end{frame}

\begin{frame}[containsverbatim]{デーモン設定（2）}

実行ファイルをディレクトリに配置

\begin{commandline}
# mkdir /usr/local/gofcgi-apid
# cp fcgi_server /usr/local/gofcgi-apid/
# chmod 755 /usr/local/gofcgi-apid/fcgi_server
\end{commandline}

serviceファイルをsystemdに登録し、繁栄

\begin{commandline}
# cp gofcgi-apid.service /lib/systemd/system/
# systemctl daemon-reload
\end{commandline}

APIサーバをsystemctlコマンドで起動/停止

\begin{commandline}
# systemctl start gofcgi-apid
# systemctl stop gofcgi-apid
\end{commandline}

\end{frame}


\section{おわりに}

\begin{frame}[containsverbatim]{おわりに}
  \begin{itemize}
  \item Debian パッケージで提供している golang を紹介しました
  \item golang でサーバアプリを作ってみました
  \item Debian で golang ライフをお楽しみください
  \end{itemize}
\end{frame}


\section{参考情報}

\begin{frame}{参考情報}
  \begin{itemize}
  \item The Go Programming Language
    \begin{itemize}
    \item \url{https://golang.org/}
    \end{itemize}
  \item Debian Go Packaging
    \begin{itemize}
    \item \url{https://go-team.pages.debian.net/packaging.html}
    \end{itemize}
  \item debianにおけるnginxの設定例
    \begin{itemize}
    \item \url{https://tokyodebian-team.pages.debian.net/pdf2018/debianmeetingresume201811-presentation-sugimoto.pdf}
    \end{itemize}
  \end{itemize}
\end{frame}

\end{document}


;;; Local Variables: ***
;;; outline-regexp: "\\([ 	]*\\\\\\(documentstyle\\|documentclass\\|emtext\\|section\\|begin{frame}\\)\\*?[ 	]*[[{]\\|[]+\\)" ***
;;; End: ***
